\section{Conclusion}
\label{sec:conclusion}

We studied how people read Trashbots across five deployment sites spanning, one in each New York City borough. Our findings show that people did not understand a public robot from its form or behavior alone. They worked out what the robot was doing by using question-driven and reasoning-based strategies to make sense of their presence. These strategies included recurring WHAT, HOW, and WHO questions as well as inference, projection, and comparison. How participants read the robots' actions was often shaped by situated knowledge of the encounter site and by personal experience brought to the interaction.

These findings motivate \emph{local legibility} as a design objective for public HRI. We identified how local legibility emerges through the interpretive process by showing how people combine cues from the robot with site-specific knowledge and personal experience to make sense of its presence. We proposed design implications for making public robot deployments more locally legible, including disclosing the responsible institution and sensing practices, relating the robot's service to existing infrastructure and activities, and accounting for the expectations that people within the community bring to the encounter. Because the same robot can be read differently across and within deployment sites, designing for local legibility also requires making the service arrangement understandable in place while recognizing that no single site-level explanation will resolve individual's interpretation.

Longer-term and fully autonomous deployments are still needed to examine how local legibility develops over time as people become familiar with a robot, experience its failures and adaptations, and form expectations for it. We hope our work can help extend discussions of legibility in public HRI beyond making a robot's individual actions understandable toward considering how its broader service role, responsibility, practices, and fit with the surrounding place become intelligible to the people who share the spaces in which it is situated.